\documentclass[12pt]{article}
\usepackage[utf8]{inputenc}
\usepackage{amsmath}
\usepackage{graphicx}
\usepackage{natbib}
\usepackage{hyperref}
\usepackage{lipsum} % For dummy text, you can remove this in your final document

\title{Impact of Climate Change on Polar Bear Populations: Analysis and Conservation Strategies}
\author{Your Name}
\date{\today}

\begin{document}

\maketitle

\begin{abstract}
Polar bears are iconic Arctic species facing significant threats due to climate change. This paper investigates the effects of climate change on polar bear populations, integrates field studies, satellite tracking data, and climate models to predict future trends, and proposes conservation strategies to mitigate these impacts.
\end{abstract}

\section{Introduction}
Polar bears (\textit{Ursus maritimus}) are apex predators adapted to Arctic sea ice habitats, crucial for their hunting and breeding. Climate change-induced sea ice loss poses a severe threat to polar bears, affecting their habitat, prey availability, and reproductive success. This paper reviews the current understanding of climate change impacts on polar bears, analyzes data from field studies and satellite tracking, and discusses conservation strategies to safeguard their populations.

\section{Data Analysis}
\subsection{Field Studies}
Field research provides critical insights into polar bear behavior, habitat use, and population dynamics under changing climatic conditions. Studies on bear health, diet shifts, and reproductive patterns offer direct observations of climate impacts \citep{stirling1999effects, amstrup2000climate}.

\subsection{Satellite Tracking Data}
Satellite telemetry data tracks polar bear movements, habitat preferences, and responses to sea ice changes. Long-term monitoring helps assess range shifts, migration patterns, and connectivity among subpopulations \citep{atwood2016rapid, lunn2016polar}.

\section{Climate Impact}
\subsection{Sea Ice Decline}
Arctic sea ice decline affects polar bears' hunting efficiency, reducing access to seal prey and increasing energetic demands during longer fasting periods \citep{serreze2007arctic, obbard2010polar}.

\subsection{Food Web Disruption}
Changes in sea ice alter Arctic food webs, impacting prey availability for polar bears and exacerbating competition for resources with other predators \citep{whiteman2013climate, ramsay2016polar}.

\section{Population Trends}
\subsection{Population Declines}
Recent studies indicate declining polar bear populations, with reduced cub survival rates and declining body condition linked to diminishing sea ice habitats \citep{regehr2007polar, laidre2008quantifying}.

\subsection{Regional Variations}
Regional variations in sea ice loss influence population dynamics differently across Arctic regions, with some subpopulations facing more acute threats than others \citep{obbard2016polar, wiig2015polar}.

\section{Conservation Strategies}
\subsection{Mitigating Climate Impacts}
Conservation efforts focus on mitigating climate change through global emissions reductions, advocating for Arctic conservation policies, and promoting sustainable development practices \citep{durner2009predicting, amstrup2010future}.

\subsection{Habitat Protection}
Protecting critical polar bear habitats, establishing marine protected areas, and implementing adaptive management strategies are crucial for ensuring resilient populations \citep{joint2013polar, dabbs2015polar}.

\section{Conclusion}
Climate change poses existential threats to polar bears, requiring urgent conservation actions to safeguard their future. Integrating field studies, satellite data, and climate models provides a robust foundation for understanding impacts and formulating effective conservation strategies in a changing Arctic environment.

\section*{References}
\bibliographystyle{plainnat}
\bibliography{prompt6_35}

\end{document}
